\documentclass[11pt]{article}
\usepackage[T1]{fontenc}
\usepackage[utf8]{inputenc}
\usepackage{lmodern,microtype}
\usepackage[margin=1in]{geometry}
\usepackage{amsmath,amssymb,amsthm,mathtools}
\usepackage{booktabs,longtable,array}
\usepackage{xcolor,listings}
\usepackage{hyperref,fancyhdr}
\hypersetup{colorlinks=true,linkcolor=blue!45!black,citecolor=blue!45!black,urlcolor=blue!45!black,
 pdftitle={Decomposing Algebraic Roots into Low-Degree Sums and Products},
 pdfauthor={OpenAI}}
\definecolor{codebg}{gray}{0.97}
\lstdefinestyle{wl}{language=Mathematica,basicstyle=\ttfamily\footnotesize,
 backgroundcolor=\color{codebg},frame=single,rulecolor=\color{black!15},
 breaklines=true,breakatwhitespace=false,columns=fullflexible,keepspaces=true,
 showstringspaces=false,tabsize=2,numbers=none,upquote=true,
 keywordstyle=\color{blue!50!black},commentstyle=\color{black!55},
 stringstyle=\color{black!75},aboveskip=10pt,belowskip=10pt}
\lstset{style=wl}
\newcommand{\Q}{\mathbb Q}
\newcommand{\C}{\mathbb C}
\newcommand{\Qbar}{\overline{\mathbb Q}}
\newcommand{\Gal}{\operatorname{Gal}}
\newcommand{\Tr}{\operatorname{Tr}}
\newcommand{\Norm}{\operatorname{N}}
\newcommand{\Res}{\operatorname{Res}}
\newcommand{\Span}{\operatorname{span}_{\Q}}
\newcommand{\adeg}{\operatorname{adeg}}
\newcommand{\code}[1]{\texttt{\detokenize{#1}}}
\newtheorem{theorem}{Theorem}[section]
\newtheorem{proposition}[theorem]{Proposition}
\newtheorem{lemma}[theorem]{Lemma}
\newtheorem{corollary}[theorem]{Corollary}
\theoremstyle{definition}\newtheorem{definition}[theorem]{Definition}
\theoremstyle{remark}\newtheorem{remark}[theorem]{Remark}
\pagestyle{fancy}\fancyhf{}
\fancyhead[L]{\small Low-degree decompositions of algebraic numbers}
\fancyhead[R]{\small Exact algorithms and certificates}
\fancyfoot[C]{\thepage}
\setlength{\headheight}{14pt}
\setlength{\emergencystretch}{3em}
\title{Decomposing Algebraic \texttt{Root}s into\\[2pt]
Low-Degree Sums and Products\\[7pt]
\large Exact Mathematica algorithms, global lower bounds,\\
and the scope of minimality}
\author{A technical article with reference implementation}
\date{6 September 2026}

\begin{document}
\maketitle
\begin{abstract}
An algebraic number is easy to combine with other algebraic numbers, but reversing
that process while minimizing the largest degree requires more than polynomial
factorization. This article gives exact solutions to both degree-nine examples
in the question and proves that their largest component degree, three, is globally
optimal, even when arbitrarily many components and arbitrary algebraic extensions
are allowed. It then develops three complementary algorithms: resultant-based
candidate generation with exact branch certification; a complete enumeration of
embedded subfields followed by rational linear algebra; and a normal-closure
reduction that gives a finite algorithm for the unrestricted additive problem.
For multiplication, an exact intersection criterion solves the two-factor problem
in any specified number field. General finite products are also addressed through
bounded exhaustive search, but are not confused with this two-factor problem.
A worked example proves that restricting summands to the field generated by the
input can change the answer from nine to six. The accompanying Wolfram Language
package, tests, and independent exact verification scripts make the distinctions
between a witness, a bounded negative result, and a global minimum explicit.
\end{abstract}

\paragraph{What is and is not delivered.}
Both requested examples are completely solved, with global optimality proofs.
The additive theory is complete without a restriction on the number or location
of summands. The supplied multiplicative field algorithm is complete for
\emph{at most two factors in a specified field}; its result is globally optimal
whenever it meets the global lower bound. A universally terminating optimizer for
arbitrary-length products in arbitrary algebraic extensions is \emph{not} claimed
by the implementation. The bounded product searches have explicitly bounded
scopes. These are mathematical distinctions, not merely implementation details.

\clearpage
\tableofcontents
\newpage

\section{The problem and its conventions}

Write
\[
 \adeg(\beta)=[\Q(\beta):\Q].
\]
For an exact algebraic number $\alpha$, define
\begin{align}
 D_+(\alpha)&=\min_{\alpha=\beta_1+\cdots+\beta_r}
                  \max_i\adeg(\beta_i),\\
 D_\times(\alpha)&=\min_{\alpha=\beta_1\cdots\beta_r}
                  \max_i\adeg(\beta_i),\qquad \alpha\ne0.
\end{align}
Here $r$ is any positive finite integer and the components lie in $\Qbar$.
For zero we set $D_\times(0)=1$, witnessed by the single rational component zero.
Rational components are allowed and have degree one. Requiring at least two
components makes no essential difference: append zero to a sum or one to a
product. Negative signs and inverses do not change absolute degree.

The coefficient field is crucial. Every displayed polynomial defining a component
must have \emph{rational} coefficients, or integer coefficients after clearing
denominators. Allowing arbitrary algebraic coefficients would trivialize the
problem: $X-\alpha$ has formal degree one over $\Q(\alpha)$ but does not make
$\alpha$ a rational number. Likewise, a nested radical of small relative degree
need not have small absolute degree.

The input polynomial may be reducible or nonminimal. The degree to optimize is
therefore obtained from \code{MinimalPolynomial}, not by inspecting the syntax of
the first argument of an arbitrary \code{Root}. Mathematica supports exact
algebraic arithmetic through \code{Root}, \code{RootReduce}, and
\code{MinimalPolynomial} \cite{root,rootreduce,minimalpolynomial}.

There are three different scopes that must not be conflated:
\begin{center}
\begin{tabular}{p{0.23\linewidth}p{0.69\linewidth}}
\toprule
Scope & Meaning\\\midrule
Unrestricted & Components may lie anywhere in $\Qbar$; their number is unbounded.\\
Ambient-field & Every component must belong to a specified number field $K$.\\
Two-component & Only a sum or product of at most two components is considered.\\
\bottomrule
\end{tabular}
\end{center}
A search with bounded coefficient height introduces a fourth restriction.
An algorithm should report these restrictions with its result.

\section{The two requested examples: exact answers and optimality}
\label{sec:examples}
Set
\[
 p(X)=X^3+X+1,\qquad q(X)=X^3-X+1,
\]
and let $a,b$ be their respective real roots:
\begin{lstlisting}
a = Root[1 + # + #^3 &, 1];
b = Root[1 - # + #^3 &, 1];
\end{lstlisting}
The two polynomials in the question, written in descending order, are
\begin{align}
 P(X)&=X^9+2X^7-3X^6+X^5-X^4+3X^3-X-1,\label{eq:P}\\
 S(X)&=X^9+6X^6+3X^5-15X^3+24X^2-4X+8.\label{eq:S}
\end{align}
The answer is exactly
\[
 \boxed{\alpha_\times=ab,\qquad \alpha_+=a+b.}
\]
The following is an executable verification, not a numerical comparison:
\begin{lstlisting}
alphaProduct = Root[
  -1 - # + 3 #^3 - #^4 + #^5 - 3 #^6 + 2 #^7 + #^9 &, 1];
alphaSum = Root[
  8 - 4 # + 24 #^2 - 15 #^3 + 3 #^5 + 6 #^6 + #^9 &, 1];

{RootReduce[a b - alphaProduct],
 RootReduce[a + b - alphaSum]}
(* {0, 0} *)
\end{lstlisting}
For orientation only,
\[
 a\approx-0.682327803828019,\quad
 b\approx-1.324717957244746,
\]
so
\[
 \alpha_\times\approx0.903891894458348,\qquad
 \alpha_+\approx-2.007045761072765.
\]
Both $P$ and $S$ have exactly one real root. These root counts, their
irreducibility, and the resultant identities below were independently checked
with exact SymPy calculations in the supplied verification scripts.

\subsection{Why the degrees really are nine}
There is also a structural proof of irreducibility. The cubics $p,q$ are
irreducible by the rational-root test, and their discriminants are respectively
$-31$ and $-23$. Their splitting fields $L_p,L_q$ have Galois groups $S_3$.
The only nontrivial proper Galois subfield of an $S_3$-extension is its quadratic
discriminant field. Those quadratic fields are different here; consequently
\[
 L_p\cap L_q=\Q,\qquad
 \Gal(L_pL_q/\Q)\simeq S_3\times S_3.
\]
These standard Galois-correspondence facts may be found in \cite{milne}.

Let $a_1,a_2,a_3$ and $b_1,b_2,b_3$ be the roots. All nine pairs occur under
the Galois action. If $a_i+b_j=a_k+b_\ell$, then
$a_i-a_k=b_\ell-b_j$ belongs to $L_p\cap L_q=\Q$.
Two conjugates cannot differ by a nonzero rational: iterate a finite-order
Galois automorphism taking one to the other. Thus $i=k,j=\ell$.

For products, a repeated value would give $a_i/a_k=b_\ell/b_j\in\Q$.
A rational ratio of nonzero conjugates is $1$ or $-1$, again by iterating an
automorphism. Neither cubic has two opposite roots, since
$p(X)+p(-X)=q(X)+q(-X)=2$. Thus the nine products are also distinct.
Both Galois orbits therefore have size nine, proving
\[
 \adeg(\alpha_+)=\adeg(\alpha_\times)=9.
\]

\subsection{A lower bound that applies to arbitrarily many components}
\begin{theorem}[Prime-divisor lower bound]\label{thm:prime}
Let $n=\adeg(\alpha)>1$, and let $\ell$ be the largest prime divisor of $n$.
Then
\[
 D_+(\alpha)\ge\ell,\qquad D_\times(\alpha)\ge\ell.
\]
The same obstruction applies to any rational expression in finitely many
algebraic numbers of degree less than $\ell$.
\end{theorem}
\begin{proof}
Suppose all components have degree at most $d$. The splitting field of a
component of degree $m\le d$ has Galois group embedded in $S_m$, so every
prime dividing its degree is at most $d$. The Galois group of the compositum
of these splitting fields embeds in their direct product. Its order therefore
has no prime divisor greater than $d$. The degree of any subfield, in particular
$\Q(\alpha)$, divides that order. Hence every prime divisor of $n$ is at most $d$.
\end{proof}
For $n=9$ this excludes all decompositions into linear and quadratic algebraic
numbers, regardless of their number. The cubic witnesses above attain the
bound. Therefore
\[
 \boxed{D_\times(\alpha_\times)=3,\qquad D_+(\alpha_+)=3.}
\]
In this special case the argument can be shortened: any finite collection of
quadratic numbers lies in a multiquadratic extension of degree a power of two;
its subfields cannot have degree nine.

\section{Resultants: generate candidates, then certify the chosen root}

The original Mathematica discussion contains a companion-matrix construction
by yarchik and a resultant construction by Daniel Lichtblau \cite{mse}.
These are useful ways to recover a prescribed shape of answer. The important
additional steps are to retain rational coefficient fields, allow the correct
polynomial divisibility condition, identify the requested conjugate, and prove
minimality separately.

\subsection{Two exact composition formulas}
For monic $f,g\in\Q[X]$ of degrees $m,n$, define
\begin{align}
 C_+(f,g)(X)&=\Res_Y\bigl(f(Y),g(X-Y)\bigr),\label{eq:sumres}\\
 C_\times(f,g)(X)&=\Res_Y\bigl(f(Y),Y^n g(X/Y)\bigr).\label{eq:prodres}
\end{align}
If $f$ has roots $u_i$ and $g$ has roots $v_j$, these are
\[
 C_+(f,g)=\prod_{i,j}(X-u_i-v_j),\qquad
 C_\times(f,g)=\prod_{i,j}(X-u_iv_j).
\]
These identities follow immediately by evaluating the second polynomial at
all roots of the first. The second expression is a polynomial: if
$g(X)=\sum_{j=0}^{n}g_jX^j$, use
\[
 Y^n g(X/Y)=\sum_{j=0}^{n}g_jX^jY^{n-j}
\]
rather than performing a rational substitution followed by unsafe cancellation.
This formula also handles zero roots correctly. Mathematica's
\code{Resultant} implements the required elimination \cite{resultant}.

For the two cubics,
\[
 C_\times(p,q)=P,\qquad C_+(p,q)=S.
\]
Using the package in the archive:
\begin{lstlisting}
Get["RootDecomposition.wl"];
ComposedPolynomial[x^3 + x + 1, x^3 - x + 1, x, "Product"]
ComposedPolynomial[x^3 + x + 1, x^3 - x + 1, x, "Sum"]
\end{lstlisting}
The implementation normalizes the input polynomials to monic form. The
resultants then already have leading coefficient one. There is no need for
an unexplained sign correction to a characteristic polynomial.

\subsection{Why polynomial equality is too restrictive in general}
Let $M_\alpha$ be the minimal polynomial of the input. A necessary and
sufficient condition for \emph{some} pair of roots of $f,g$ to produce
$\alpha$ is
\[
 M_\alpha\mid C_\circ(f,g),\qquad \circ\in\{+,\times\}.
\]
When $mn=\adeg(\alpha)$, monicity makes divisibility equivalent to equality.
Without that degree equality the composed polynomial may be reducible or
have repeated roots.

For example, take $f=X^2-2$ and $g=X^2-8$. One sum of their roots is $\sqrt2$,
using $-\sqrt2+\sqrt8=\sqrt2$. The composed polynomial is
\[
 (X^2-2)(X^2-18),
\]
not the minimal polynomial $X^2-2$. Searching only by equating a composed
polynomial to the input would miss this valid pair.

Nor is $\adeg(\alpha)$ generally a divisor of the product of component degrees.
If $c=\sqrt[3]{2}$ and $\omega=e^{2\pi i/3}$, then
\[
 c+2\omega c=i\sqrt3\,\sqrt[3]{2}
\]
has minimal polynomial $X^6+108$, although each summand has degree three.
The degree of the compositum is six, not nine. The safe universal statement
is an \emph{inequality} $\adeg(\alpha)\le\prod_i\adeg(\beta_i)$; stronger
claims need additional hypotheses.

\subsection{Coefficient matching is a proposal mechanism}
For proposed degrees $m,n$, introduce generic monic rational polynomials
$f,g$ and set the coefficients of
\[
 C_\circ(f,g)\bmod M_\alpha
\]
to zero. If $mn=\deg M_\alpha$, one may instead equate all coefficients of
$C_\circ(f,g)$ and $M_\alpha$.
This produces a polynomial system in the unknown coefficients. Solving it
over $\C$ or $\mathbb R$ is not the same as solving it over $\Q$.
An algebraic solution for the coefficients can violate the original objective.
A parameterized answer from \code{Solve} also needs an actual admissible
specialization and subsequent exact checks \cite{solve}.

Translations $u\mapsto u+r$, $v\mapsto v-r$ with $r\in\Q$ preserve a sum.
Rescalings $u\mapsto ru$, $v\mapsto v/r$ with $r\in\Q^\times$ preserve a
product. They can help reduce a symbolic search, but multiplicative gauges
require care: setting the constant coefficient of a monic degree-$m$
polynomial to one by rational rescaling requires a suitable rational $m$th
root. That is not always possible. The package avoids assuming such a gauge.

\subsection{Certifying the actual conjugate}
A minimal polynomial does not specify which conjugate was intended. An exact
identity is checked with
\begin{lstlisting}
RootReduce[candidate - alpha] === 0
\end{lstlisting}
not with numerical closeness. The package enumerates roots of the proposed
polynomials and tests each appropriate pair:
\begin{lstlisting}
productAnswer = PairFromPolynomials[
  alphaProduct, x^3 + x + 1, x^3 - x + 1, x, "Product"];
sumAnswer = PairFromPolynomials[
  alphaSum, x^3 + x + 1, x^3 - x + 1, x, "Sum"];

productAnswer["Components"]
(* {Root[1 + # + #^3 &, 1], Root[1 - # + #^3 &, 1]} *)
productAnswer["GlobalMinimumCertified"]
(* True *)
\end{lstlisting}
The displayed comments indicate mathematically expected results; the native
Wolfram tests supplied with the archive were not executed in this environment.

A cheaper candidate search guesses only one polynomial. For each of its roots
$u$, compute $v=\alpha-u$ or $v=\alpha/u$ and measure its actual degree. This
avoids guessing the second polynomial altogether:
\begin{lstlisting}
PairFromOnePolynomial[alphaProduct, x^3 + x + 1, x, "Product", 3]
BoundedTwoRootSearch[alphaProduct, "Product", 3, 1]
\end{lstlisting}
The second call enumerates primitive integer first-component polynomials of
degree at most three and coefficient height at most one. The polynomial
$X^3+X+1$ lies in that search. Its complementary factor is not required to
satisfy the same height bound. A failed bounded search proves only the
corresponding bounded two-component statement, not unrestricted impossibility.

\section{A finite, exact enumeration of all embedded subfields}
\label{sec:subfields}
The next method requires no guess about the factor polynomials. It is not
intended to replace sophisticated subfield algorithms, but it is small enough
to implement directly in Wolfram Language and to prove from first principles.
More efficient approaches are discussed in \cite{subfields,pari}.

Let $K=\Q(\theta)$ have degree $N$, with monic minimal polynomial $M(T)$.
Factor it over $K$:
\[
 M(T)=(T-\theta)g_1(T)\cdots g_s(T),
\]
where all factors are monic and irreducible in $K[T]$. Characteristic zero
ensures squarefreeness. Mathematica provides such factorization through
\code{FactorList[..., Extension -> theta]} \cite{factor}.

\begin{theorem}[Subfields from coefficient fields of divisors]\label{thm:subfields}
Every intermediate field $\Q\subseteq E\subseteq K$ is generated over $\Q$ by
the coefficients of a divisor
\[
 h(T)=(T-\theta)\prod_{j\in J}g_j(T)
\]
for some subset $J\subseteq\{1,\ldots,s\}$. If $h$ is the minimal polynomial
of $\theta$ over $E$, then
\[
 [E:\Q]=N/\deg h.
\]
Conversely, if the coefficient field of such an $h$ has this degree, $h$ is
its minimal polynomial for $\theta$.
\end{theorem}
\begin{proof}
Take the minimal polynomial $h_E$ of $\theta$ over $E$. It divides $M$ in
$E[T]$, hence in $K[T]$, and contains the factor $T-\theta$. Its coefficients
therefore arise from a divisor of the asserted form.
Let $F$ be the field generated by those coefficients. Then $F\subseteq E$,
while $h_E(\theta)=0$ gives
\[
 [K:F]\le\deg h_E=[K:E].
\]
The reverse inequality follows from $F\subseteq E$, so $F=E$.
Conversely, if $F$ is the coefficient field of $h$, the minimal polynomial of
$\theta$ over $F$ divides $h$. Equality of their degrees gives equality of
these monic polynomials.
\end{proof}

Thus the following finite procedure returns every embedded subfield. Enumerate
the $2^s$ subsets $J$. Form $h$. Skip it unless $\deg h$ divides $N$. Compute
the rational algebra generated by its coefficients. Retain it when its dimension
is $N/\deg h$, and remove duplicate embedded fields.

\subsection{Computing a coefficient field by rational linear algebra}
Represent an element of $K$ by its coordinate row in
\[
 1,\theta,\ldots,\theta^{N-1}.
\]
Products of coordinate rows are computed by multiplying the corresponding
polynomials and reducing modulo $M$. Starting with the row for one, repeatedly
adjoin products of the current basis rows with the coefficient generators,
then row-reduce. The dimension can increase at most $N-1$ times.

The final space is a finite-dimensional $\Q$-subalgebra of the field $K$.
It is itself a field: multiplication by a nonzero element is injective on
that finite-dimensional algebra, hence surjective, so its inverse belongs to
the algebra. This proves the closure computation correct.

Canonical reduced row spaces identify \emph{embedded} subfields, not just
isomorphism classes. Two isomorphic subfields in different positions in $K$
can contribute differently to the desired decomposition and must not be merged.

The package first replaces a nonintegral generator by an integer multiple
that is integral. This is an exact change of generator, not a change of field.
It matters because \code{ToNumberField[b,theta]} guarantees the indicated
\code{AlgebraicNumber[theta,...]} basis when $\theta$ is integral; it can
otherwise normalize the generator \cite{tonumberfield,algebraicnumber}.
The code verifies the generator and the rational coordinates rather than
silently assuming a basis.

\subsection{What happens for the two degree-nine examples}
For both inputs the irreducible factor degrees of $M$ over $K$ are
\[
 1,2,2,4.
\]
Only eight divisors containing the distinguished linear factor need be
examined. The embedded subfield degrees returned are
\[
 \boxed{1,3,3,9.}
\]
The two degree-three subfields are $\Q(a)$ and $\Q(b)$. They are discovered
from the degree-nine input alone; neither cubic is supplied to the algorithm.
\begin{lstlisting}
AllSubfieldBases[alphaProduct]["Degrees"]
(* {1, 3, 3, 9} *)
AllSubfieldBases[alphaSum]["Degrees"]
(* {1, 3, 3, 9} *)
\end{lstlisting}
The pure subset enumeration is exponential in the number of factors. In a
Galois ambient field of degree $N$, the minimal polynomial of a primitive
element splits into $N$ linear factors, making the count $2^{N-1}$.
The default \code{"MaxDivisors"} guard returns an explicit resource failure;
it never converts an incomplete enumeration into a negative mathematical answer.

\section{All finite sums inside a specified number field}
For a finite extension $K/\Q$ and a degree bound $d$, define
\[
 V_d(K)=\sum_{\substack{\Q\subseteq E\subseteq K\\{}[E:\Q]\le d}} E,
\]
where the sum is a sum of $\Q$-vector subspaces, not a compositum of fields.

\begin{theorem}[Additive membership criterion]\label{thm:additive}
An element $\alpha\in K$ is a finite sum of elements of absolute degree at
most $d$ in $K$ if and only if $\alpha\in V_d(K)$.
\end{theorem}
\begin{proof}
A permitted summand $\beta$ belongs to the subfield $\Q(\beta)$ of degree at
most $d$, proving necessity. Conversely, each element of a degree-at-most-$d$
subfield has absolute degree at most $d$. A vector-space expression in the
sum of these fields is therefore a permitted additive decomposition.
\end{proof}

Once subfield bases are known, stack their coordinate rows into a matrix $B$.
The decision is the exact consistency test
\[
 B^{\mathsf T}c=v_\alpha.
\]
Solving this rational system gives the summands. All scalar coefficients are
rational and can be absorbed into the summands without increasing their degrees.
Grouping contributions from the same subfield cannot increase the maximum degree.
Only the distinct subfield degrees need be tried, in increasing order.

This handles \emph{any finite number of summands}; it is not a binary search.
Indeed a basis selected from the union of all permitted field bases uses at most
$[K:\Q]$ elements, so that many summands always suffice whenever membership holds.

The package also centers each nonrational summand by subtracting its mean trace
\[
 \mu(\beta)=\frac{\Tr_{\Q(\beta)/\Q}(\beta)}{\adeg(\beta)}.
\]
It collects the removed rational quantities into one rational summand. This
preserves the identity and the degree bound. For the requested sum example it
removes arbitrary rational translations introduced by linear algebra and recovers
the two depressed cubics.
\begin{lstlisting}
s = MinSumInField[alphaSum];
s["MaxDegree"]                    (* 3 *)
s["ExactIdentityVerified"]        (* True *)
s["GlobalMinimumCertified"]       (* True *)
\end{lstlisting}
The last statement is justified by Theorem~\ref{thm:prime}, not merely by the
fact that the internal search was exhaustive.

\section{Two factors inside a specified field: an intersection test}
\label{sec:product}
Multiplication is not a linear operation, but its two-factor version has an
especially simple linear reformulation.

\begin{theorem}[Two-factor intersection criterion]\label{thm:intersection}
Let $E,F\subseteq K$ be subfields and $\alpha\in K^\times$. Then
\[
 \alpha\in E^\times F^\times
 \quad\Longleftrightarrow\quad E\cap\alpha F\ne\{0\}.
\]
Given a nonzero element $u=\alpha v$ of that intersection, with $u\in E$ and
$v\in F$, a factorization is $\alpha=u\,v^{-1}$.
\end{theorem}
\begin{proof}
If $\alpha=bc$ with $b\in E^\times,c\in F^\times$, then
$b=\alpha c^{-1}$ lies in the intersection. Conversely, a nonzero
$u=\alpha v$ gives $\alpha=u/v$, and $v^{-1}$ belongs to $F$.
\end{proof}

Let $U$ and $V$ be row bases for $E$ and $F$. Multiply every row of $V$ by
$\alpha$, producing the row matrix $W$. A nonzero rational null vector of
\[
 \begin{bmatrix}U^{\mathsf T}&-W^{\mathsf T}\end{bmatrix}
\]
gives coefficients for $u$ and $v$. Since $U,V$ are bases and $\alpha\ne0$,
a nonzero null vector cannot give $u=v=0$. Thus a nontrivial nullspace is
exactly the desired condition. No norm equation, logarithm, or rational-point
solver is needed for this two-field problem.

Enumerate all pairs of embedded subfields, sorted by
$\max([E:\Q],[F:\Q])$, and apply this test. The first success minimizes the
largest degree among all two-factor decompositions \emph{inside the ambient
field}. This remains true when a recovered factor generates a smaller field:
that smaller subfield also occurs in the enumeration.
\begin{lstlisting}
r = MinTwoFactorInField[alphaProduct];
r["MaxDegree"]                    (* 3 *)
r["ComponentDegrees"]             (* {3, 3} *)
r["ExactIdentityVerified"]        (* True *)
r["GlobalMinimumCertified"]       (* True *)
\end{lstlisting}
A basis-dependent answer may be a rationally rescaled version of $a,b$.
For example, the independent implementation produced $-2a/3$ and $-3b/2$,
whose minimal polynomials are
\[
 X^3+\frac49X-\frac8{27},\qquad
 X^3-\frac94X-\frac{27}{8}.
\]
Their product is exactly $ab$. The degree objective is unchanged by rational
rescaling. The earlier \code{PairFromPolynomials} call gives the particular
small-height representatives stated in the question.

\subsection{Why two factors do not settle arbitrary finite products}
Consider
\[
 \gamma=(1+\sqrt2)(1+\sqrt3)(1+\sqrt5).
\]
It has degree eight and is a product of three quadratic numbers. No product
of only two quadratic numbers can have degree eight, since their compositum
has degree at most four. Thus optimizing over two components can give a
strictly worse answer than optimizing over all finite products.

In a fixed ambient field, the arbitrary-product membership problem is
\[
 \alpha\in H_d(K):=
 \left\langle E^\times:\ E\subseteq K,\ [E:\Q]\le d\right\rangle.
\]
There are finitely many such subfields, and factors from the same subfield can
be combined. Nevertheless, testing a product of many multiplicative groups
is not the same as testing one pairwise intersection. Greedy recursive splitting
is a useful search strategy but does not by itself prove the global optimum.

\section{A complete reduction for the unrestricted additive problem}
\label{sec:global}
The preceding additive algorithm becomes globally complete when its ambient
field is the normal closure of the input field. The point is not just that
this field contains many conjugates: normalized trace preserves the component
degree bound.

\begin{theorem}[Normal-closure descent for sums]\label{thm:descent}
Let $L/\Q$ be finite Galois and contain $\alpha$. For every $d\ge1$,
\[
 \alpha\text{ is a finite sum of numbers of degree at most }d
 \quad\Longleftrightarrow\quad \alpha\in V_d(L).
\]
Thus the unrestricted additive minimum is attained by a decomposition entirely
inside the normal closure of $\Q(\alpha)$.
\end{theorem}
\begin{proof}
Suppose $\alpha=\sum_i\beta_i$ with $\adeg(\beta_i)\le d$.
Choose a finite Galois extension $M/\Q$ containing $L$ and all the $\beta_i$.
Set $G=\Gal(M/\Q)$ and $N=\Gal(M/L)$. Since $L/\Q$ is Galois, $N$ is normal
in $G$. Define the rational-linear averaging operator
\[
 \pi(z)=\frac1{|N|}\sum_{\sigma\in N}\sigma(z).
\]
It fixes $L$, maps $M$ to $L$, and commutes with the action of $G$.
Therefore
\[
 \alpha=\sum_i\pi(\beta_i).
\]
The map from the $G$-orbit of $\beta_i$ to that of $\pi(\beta_i)$ is
surjective and equivariant: $g\beta_i\mapsto\pi(g\beta_i)=g\pi(\beta_i)$.
Orbit size is absolute algebraic degree, so
\[
 \adeg\bigl(\pi(\beta_i)\bigr)\le\adeg(\beta_i)\le d.
\]
Zero summands may be omitted. This proves the forward direction. The reverse
direction is Theorem~\ref{thm:additive} applied to $L$.
\end{proof}

This proves a genuinely finite algorithm: construct the normal closure, enumerate
its subfields, and use rational linear algebra at successively larger degree
thresholds. There is no coefficient-height bound and no unbounded search over
numbers of summands. Termination here is a mathematical statement, not a claim
of acceptable running time on every input.

\subsection{Implementation and practical limits}
To construct the normal closure, list every exact root of the input's minimal
polynomial and put them in their smallest common number field. The package uses
\code{ToNumberField[roots, All]}; the \code{All} specification matters because
the default common-field construction need not use the smallest common extension
\cite{tonumberfield}.

The public entry point is
\begin{lstlisting}
MinSumGlobally[alphaSum]
\end{lstlisting}
It first searches inside $\Q(\alpha)$. If that witness meets the prime-divisor
lower bound, global optimality is already established and no normal closure is
constructed. This is what happens for the requested sum example.

Otherwise it constructs the normal closure and invokes the complete field
procedure. The package has explicit degree and divisor-count guards. Removing
those guards gives the finite exhaustive algorithm described by the proof:
\begin{lstlisting}
MinSumGlobally[alpha,
  "MaxDivisors" -> Infinity, "MaxNormalDegree" -> Infinity]
\end{lstlisting}
This is not a performance recommendation. The normal closure can have degree
as large as $n!$, and the elementary subfield enumeration can be much worse
than the dedicated algorithms in \cite{subfields,pari}. In particular, the
degree guard is checked after the common-field construction; an outer
\code{TimeConstrained} may be needed to bound that construction as well.
A timeout or resource failure proves no nonexistence assertion.

\section{A concrete warning: internal degree nine, global degree six}
\label{sec:counterexample}
The same product input $\alpha_\times=ab$ gives a particularly instructive
counterexample to silently restricting summands to $\Q(\alpha)$.

\subsection{The internal answer is nine}
The subfields of $K=\Q(ab)=\Q(a,b)$ have degrees $1,3,3,9$.
The sum of its proper subfields is
\[
 \Q(a)+\Q(b)=\Span\{1,a,a^2,b,b^2\}.
\]
The fields have compositum degree nine, so
\[
 \{a^i b^j:0\le i,j\le2\}
\]
is a basis of $K$. In particular $ab$ does not lie in that five-dimensional
sum. The internal additive minimum is therefore nine. The exact linear-algebra
verification script obtains the same result.

\subsection{An explicit global degree-six decomposition}
Let $a_1=a,b_1=b$ and label each remaining pair by complex conjugation.
Define
\begin{align*}
 t_1&=a_1b_1+a_2b_2+a_3b_3,\\
 t_2&=a_1b_1+a_2b_3+a_3b_2.
\end{align*}
Since each cubic has zero trace,
\[
 t_1+t_2=2a_1b_1+(a_2+a_3)(b_2+b_3)=3a_1b_1.
\]
The six bijective matching sums
$\sum_i a_i b_{\sigma(i)}$, $\sigma\in S_3$, form a Galois-stable set.
Their polynomial is
\[
 R(T)=T^6+6T^4-27T^3+9T^2-81T+4.
\]
One exact elimination certificate is
\begin{multline}
 \Res_U\!\left(U^3+U+1,
  \Res_V\!\left(V^3-V+1,
    T^2-3UVT+3U^2-3V^2+4\right)\right)\\
 =R(T)^3.
\end{multline}
Indeed $t_1+t_2=3ab$ and $t_1t_2=3a^2-3b^2+4$.
The polynomial $R$ is irreducible and has exactly two real roots, namely
$t_1,t_2$. Consequently
\begin{lstlisting}
r1 = Root[4 - 81 # + 9 #^2 - 27 #^3 + 6 #^4 + #^6 &, 1];
r2 = Root[4 - 81 # + 9 #^2 - 27 #^3 + 6 #^4 + #^6 &, 2];
RootReduce[(r1 + r2)/3 - alphaProduct]
(* 0 *)
\end{lstlisting}
The summands $r_1/3,r_2/3$ each have degree six. They can individually be
converted to rational-polynomial \code{Root} objects by \code{RootReduce}.

\subsection{Why six is globally minimal}
This stronger lower bound has a short group-theoretic proof.
Work in $L=L_pL_q$ with $G=S_3\times S_3$.
Let $A_1,A_2$ be the normal cyclic subgroups of order three in the two factors,
and let $e_i$ be averaging over $A_i$. Put
\[
 T=(1-e_1)(1-e_2).
\]
The zero traces of $a,b$ give $T(ab)=ab\ne0$.

Now let $\beta\in L$ have degree at most five and let $H$ be its stabilizer.
Its index divides $36$, so it is $1,2,3$, or $4$. If the index is $1,2$, or
$4$, then $H$ contains the unique Sylow three-subgroup $A_1A_2$, so $T\beta=0$.
If the index is three, $H$ has order twelve and $H\cap(A_1A_2)$ has order
three. Moreover $H$ maps onto
$G/(A_1A_2)\simeq C_2\times C_2$. Its intersection with $A_1A_2$ must
therefore be invariant under the two independent inversions on
$A_1A_2\simeq\mathbb F_3^2$. The only invariant lines are the two coordinate
axes. Hence $H$ contains $A_1$ or $A_2$, again giving $T\beta=0$.

Thus $T$ annihilates every element of $L$ of degree at most five, and therefore
every sum of such elements. It does not annihilate $ab$. Theorem~\ref{thm:descent}
extends the obstruction to summands outside $L$. We have proved
\[
 \boxed{D_+(ab)=6,\qquad D_+^{\,\Q(ab)}(ab)=9.}
\]
The normal-closure step is consequently essential in general, not merely
an overly cautious implementation choice.

\section{What can be asserted for unrestricted products?}
The exact two-factor criterion and the global prime-divisor bound already
solve the product example completely. They should not be promoted, without
another theorem, to a general arbitrary-product algorithm.

There is a tempting but invalid shortcut using norms. If a factorization is
defined in a finite Galois extension $M/\Q$, and $L/\Q$ is a Galois
subextension containing $\alpha$, then
\[
 \alpha^{[M:L]}=\prod_i\Norm_{M/L}(\beta_i).
\]
Each norm belongs to $L$. Moreover, by the same equivariance argument as for
trace, its absolute degree is at most that of $\beta_i$. But this proves a
statement about a \emph{power} of $\alpha$, not about $\alpha$ itself.
Extracting roots can increase absolute degrees and introduces root-of-unity
ambiguities. For example, $\sqrt[3]{2}^{\,3}=2$ is rational, but
$\sqrt[3]{2}$ is not a product of quadratic numbers by
Theorem~\ref{thm:prime}. Merely testing membership of a power in a low-degree
multiplicative subgroup is therefore insufficient.

\subsection{An exact bounded search for any number of factors}
Fix degree, coefficient-height, and component-count bounds $d,H,R$.
Enumerate every primitive integer polynomial of degree at most $d$, positive
leading coefficient, and coefficients of absolute value at most $H$.
Keep the irreducible polynomials and enumerate all their exact roots. This
is a finite catalog containing precisely the possible components within
that polynomial-height convention. Enumerate tuples of length at most $R$
and check their products by \code{RootReduce}. The same construction works
for sums. A finite search over this entire box can determine its minimum
largest degree as well, by examining bounds in increasing order.

The package includes \code{FiniteCatalogSearch} for this last stage. Repetitions
are allowed, and all successes are exact certificates. Expanding $H$ and $R$
in a fair order eventually finds any existing factorization at a specified
bound $d$. This is a positive semidecision procedure: failure at all stages
cannot be recognized simply by waiting for the search to finish. In particular,
one must not search $d=1,2,3,\ldots$ sequentially with an unbounded search at
each $d$ and claim to obtain an optimizer.

For production work, a useful division of labor is to run field-based searches,
small-height searches, and educated resultant ansatzes as proposal mechanisms,
while a separate exact checker verifies identities, degrees, and the scope of
any optimality assertion. Ideal, unit, and multiplicative-relation methods may
also be useful inside a fixed field, but require their own completeness proofs.

\section{Using the supplied implementation}

The archive contains a loadable package rather than notebook-specific cells.
A basic session is:
\begin{lstlisting}
Get["RootDecomposition.wl"];

productResult = MinTwoFactorInField[alphaProduct];
sumResult = MinSumGlobally[alphaSum];

productResult["Components"]
sumResult["Components"]
{productResult["MaxDegree"], sumResult["MaxDegree"]}
(* {3, 3} *)
\end{lstlisting}
For another ambient field, pass its generator explicitly:
\begin{lstlisting}
MinSumInField[alpha, "AmbientGenerator" -> theta]
MinTwoFactorInField[alpha, "AmbientGenerator" -> theta]
\end{lstlisting}
The code checks that the input belongs to that field. It returns associations
containing the components, their individual degrees, an inactive expression,
an exact-identity flag, a scope description, and the global prime-divisor
lower bound. \code{"GlobalMinimumCertified"} is true only when the witness
meets that bound or the complete normal-closure additive procedure establishes
the optimum. A \code{Failure} describing a resource limit is not a decomposition
and is not a proof of impossibility.

\begin{center}
\begin{tabular}{p{0.40\linewidth}p{0.52\linewidth}}
\toprule
Entry point & Guarantee when it succeeds\\\midrule
\code{PairFromPolynomials} & Exact branch-certified witness for the specified polynomials.\\
\code{BoundedTwoRootSearch} & Exact witness in a bounded first-polynomial search; no general negative certificate.\\
\code{AllSubfieldBases} & Complete list of embedded subfields, unless an explicit failure is returned.\\
\code{MinSumInField} & Minimum for all finite sums in the specified field.\\
\code{MinTwoFactorInField} & Minimum for at most two factors in the specified field.\\
\code{MinSumGlobally} & Unrestricted additive minimum, when completed.\\
\code{FiniteCatalogSearch} & Exact witness using a finite catalog and bounded component count.\\
\bottomrule
\end{tabular}
\end{center}

\subsection{Validation status}
The Wolfram computation service available during preparation returned HTTP 404
errors, and no local Wolfram kernel was installed. Consequently this article
does not claim that \code{RootDecomposition.wl} or \code{tests.wlt} was executed
in Mathematica. The Wolfram implementation was checked for balanced syntax
and reviewed against the cited documented primitives.

Independent exact implementations in SymPy \emph{were} executed. They verified:
the two degree-nine resultants, irreducibility and real-root counts; the
$1,2,2,4$ factor-degree patterns over the input fields; the complete subfield
degree lists $1,3,3,9$; the additive row-space and multiplicative nullspace
algorithms; the internal minima summarized below; and the sextic matching-sum
elimination, irreducibility, and real-root count. Additional independent checks cover the degree-eight three-factor example
and the degree-six counterexample to degree-product divisibility. The tests and logs are included rather than
being reported as native Mathematica output.
\begin{center}
\begin{tabular}{lccc}
\toprule
Input & Degree & Internal additive minimum & Internal two-factor minimum\\\midrule
$ab$ & 9 & 9 & 3\\
$a+b$ & 9 & 3 & 9\\\bottomrule
\end{tabular}
\end{center}
The second row's two-factor value is not an assertion about arbitrary-length
products outside the input field. Separately proved global results are
$D_\times(ab)=3$, $D_+(a+b)=3$, and $D_+(ab)=6$.

\subsection{Reproducibility}
Run the Wolfram test file with
\begin{lstlisting}
TestReport["tests.wlt"]
\end{lstlisting}
from the extracted directory. The independent checks can be rerun with
\begin{lstlisting}[language={}]
python verify.py
python verify_algorithms.py
\end{lstlisting}
They require SymPy. Compile the article from the same directory using
\begin{lstlisting}[language={}]
pdflatex article.tex && pdflatex article.tex
\end{lstlisting}
The full package is typeset in the appendix from the same \code{.wl} file that
is distributed, avoiding two diverging copies of the implementation.

\section{Conclusions}
For the two examples, the implementation can discover the relevant cubic
subfields automatically and certify decompositions of largest degree three.
A prime-divisor argument then proves that no lower degree is possible, even
with arbitrarily many components.

The broader additive problem has a clean exact resolution: normalized trace
reduces it to a finite Galois field, and low-degree subfield spans turn it into
rational linear algebra. The multiplicative two-factor problem has a similarly
clean local criterion, $E\cap\alpha F\ne0$, but arbitrary-length products are
a different optimization problem. Keeping those scopes explicit is what turns
a plausible simplification routine into a mathematically reliable tool.

\begingroup\small
\begin{thebibliography}{20}
\setlength{\itemsep}{3pt}
\bibitem{mse} Vladimir Reshetnikov, ``Factor a polynomial Root into Roots of smallest possible degree,'' Mathematica Stack Exchange, question 105933 (2016), with answers by yarchik and Daniel Lichtblau (2019).
\url{https://mathematica.stackexchange.com/questions/105933/}.
\bibitem{root} Wolfram Research, \emph{Root}, Wolfram Language documentation.
\url{https://reference.wolfram.com/language/ref/Root.html}.
\bibitem{rootreduce} Wolfram Research, \emph{RootReduce}, Wolfram Language documentation.
\url{https://reference.wolfram.com/language/ref/RootReduce.html}.
\bibitem{minimalpolynomial} Wolfram Research, \emph{MinimalPolynomial}, Wolfram Language documentation.
\url{https://reference.wolfram.com/language/ref/MinimalPolynomial.html}.
\bibitem{resultant} Wolfram Research, \emph{Resultant}, Wolfram Language documentation.
\url{https://reference.wolfram.com/language/ref/Resultant.html}.
\bibitem{factor} Wolfram Research, \emph{Factor}, including algebraic extension options, Wolfram Language documentation.
\url{https://reference.wolfram.com/language/ref/Factor.html}.
\bibitem{tonumberfield} Wolfram Research, \emph{ToNumberField}, Wolfram Language documentation.
\url{https://reference.wolfram.com/language/ref/ToNumberField.html}.
\bibitem{algebraicnumber} Wolfram Research, \emph{AlgebraicNumber}, Wolfram Language documentation.
\url{https://reference.wolfram.com/language/ref/AlgebraicNumber.html}.
\bibitem{solve} Wolfram Research, \emph{Solve}, Wolfram Language documentation.
\url{https://reference.wolfram.com/language/ref/Solve.html}.
\bibitem{milne} J. S. Milne, \emph{Fields and Galois Theory}, version 5.10, September 2022. In particular Sections 3 and 5.
\url{https://www.jmilne.org/math/CourseNotes/FT.pdf}.
\bibitem{subfields} Andreas-Stephan Elsenhans and J\"urgen Kl\"uners, ``Computing subfields of number fields and applications to Galois group computations,'' preprint arXiv:1610.06837 (2016).
\url{https://arxiv.org/abs/1610.06837}.
\bibitem{pari} The PARI Group, \emph{Catalogue of GP/PARI Functions: General Number Fields}, entry \code{nfsubfields}.
\url{https://pari.math.u-bordeaux.fr/dochtml/html-stable/General_number_fields.html}.
\end{thebibliography}
\endgroup

\appendix
\newpage
\section{Complete Wolfram Language reference implementation}
\label{app:code}
The code deliberately favors an inspectable exact algorithm over asymptotically
optimal subfield enumeration. It uses rational-polynomial \code{Root} objects,
rejects inexact input, and reports search limits explicitly.
\lstinputlisting[basicstyle=\ttfamily\scriptsize]{RootDecomposition.wl}

\end{document}
